%%BLOCK after="A necessary directedness condition"
\begin{dcomment}{development}{p36:pourrait être valable.}{Total asphericity is not filteredness}
The question whether $T(A)$ forces $A$ to be filtered has a negative answer in general and a positive one for monoids and for categories with finitely many objects. Both rest on a description of $T(A)$ in which no target category appears. For a functor $H:A\to\mathrm{Set}$ write $\operatorname{El}(H)$ for its category of elements: the objects are the pairs $(a,h)$ with $h\in H(a)$, and an arrow $(a,h)\to(a',h')$ is an arrow $v:a\to a'$ with $H(v)h=h'$. The statements below were not found in the sources consulted, except where a reference is given; no priority is claimed.

\textbf{Theorem (a criterion on set-valued functors).}\ours
\textit{For a small category $A$ the following are equivalent: (i) $T(A)$; (ii) $A$ is nonempty and connected, and for every functor $H:A\to\mathrm{Set}$ every connected component of $\operatorname{El}(H)$ has the integral homology of a point; (iii) the same with ``is weakly contractible''. If they hold, every connected component $E$ of every $\operatorname{El}(H)$ satisfies $T(E)$.}

\textit{Proof.}
Let $K:A\to\mathrm{Set}$ have nonempty connected $\operatorname{El}(K)$, and let $A_K$ be $A$ with one object $x$ adjoined, $A_K(x,a)=K(a)$, no arrows into $x$ but its identity, and composition given by the action of $A$ on $K$. For the inclusion $A\to A_K$ the category $x\backslash A$ is $\operatorname{El}(K)$, and $a\backslash A$ has an initial object for $a\in A$; so the inclusion is cofinal. A chain of arrows of $A_K$ not contained in $A$ consists of $x$ (possibly repeated) followed by a chain of $\operatorname{El}(K)$. Hence $N(A_K)$ is $NA$ with the cone on $N\operatorname{El}(K)$ attached along the forgetful map, and
\[
 \widetilde H_{j+1}(BA_K)\longrightarrow\widetilde H_j(B\operatorname{El}(K))\longrightarrow\widetilde H_j(BA)
\]
is exact. Under (i) both $BA$ and $BA_K$ are weakly contractible, so $\operatorname{El}(K)$ has the homology of a point. A component of $\operatorname{El}(H)$ for an arbitrary $H$ is $\operatorname{El}(K)$ for the subfunctor $K$ of the elements lying in it. This proves (i)$\Rightarrow$(ii).

(ii)$\Rightarrow$(iii). Let $E$ be a component of $\operatorname{El}(H)$ and $L:E\to\mathrm{Set}$ a functor. Then $\operatorname{El}(L)\cong\operatorname{El}(H')$ for the functor $H'(a)=\coprod_hL(a,h)$ on $A$, the sum over the $h\in H(a)$ with $(a,h)\in E$. So $E$ again has property (ii). Let $\Pi E$ be the groupoid obtained by inverting all arrows of $E$, fix an object $e_0$, and identify $G=\operatorname{Aut}_{\Pi E}(e_0)$ with $\pi_1(BE,e_0)$. For a subgroup $\Gamma\subseteq G$ the functor $L_\Gamma(e)=\operatorname{Hom}_{\Pi E}(e_0,e)/\Gamma$ has bijective transition maps, and $\operatorname{El}(L_\Gamma)\to E$ is the connected covering with fundamental group $\Gamma$. By (ii) for $E$ its first homology $\Gamma^{\mathrm{ab}}$ vanishes. Taking $\Gamma$ cyclic gives $G=1$. A simply connected CW complex with the homology of a point is weakly contractible.

(iii)$\Rightarrow$(i). Let $f:A\to B$ be cofinal. For $b\in B$ the category $b\backslash A$ is $\operatorname{El}(a\mapsto B(b,f(a)))$; it is nonempty and connected, hence weakly contractible by (iii), and Quillen's Theorem~A makes $Bf$ a weak equivalence. The constant functor with value a point shows that $BA$ is weakly contractible. So is $BB$. The last assertion was shown in the second step.\hfill$\square$

In particular the integral-homology form of $T(A)$ (every cofinal image has the homology of a point) is equivalent to $T(A)$. With coefficients in a nonzero commutative ring $R$ the same proof, with the homology form of Theorem~A, gives: every cofinal image of $A$ has the $R$-homology of a point if and only if $A$ is nonempty and connected and the components of all $\operatorname{El}(H)$ do. For $R=\mathbf Q$ this is weaker: a finite group satisfies it.

\textbf{Proposition (exact colimits suffice).}\ours
\textit{If $A$ is nonempty and connected and $\varinjlim_A$ is exact on functors $A\to\mathrm{Ab}$, then $T(A)$.}

\textit{Proof.}
For $H:A\to\mathrm{Set}$ and $F:\operatorname{El}(H)\to\mathrm{Ab}$ put $(\Sigma_HF)(a)=\bigoplus_{h\in H(a)}F(a,h)$. The functor $\Sigma_H$ is exact and $\varinjlim_{\operatorname{El}(H)}F=\varinjlim_A\Sigma_HF$. So colimits over $\operatorname{El}(H)$ are exact, and so are colimits over a component $E$ (extend by zero). The left derived functors of $\varinjlim_E$ on the constant functor $\mathbf Z$ are the groups $H_j(BE)$; they vanish for $j>0$. Criterion (ii) applies.\hfill$\square$

\textbf{Theorem (a category with $T$ that is not filtered).}\ours
\textit{Let $\Delta'$ be the category of the ordered sets $[n]=\{0,\dots,n\}$, $n\geq0$, and strictly increasing maps. Then $T(\Delta')$ holds and $\Delta'$ is not filtered; it is not sifted either.}

\textit{Proof.}
The two maps $[0]\to[1]$ are not equalized by any injection, so $\Delta'$ is not filtered. A pair of maps $[0]\to[N]\leftarrow[0]$ is two points of $[N]$, and whether the first is smaller than, equal to or larger than the second is preserved by increasing injections; so the category of such pairs has three components and the diagonal of $\Delta'$ is not cofinal.

For exactness we show: \emph{for every $n$ there are $M$ and an integral combination $t=\sum_gt_g\,g$ of maps $g:[n]\to[M]$ of $\Delta'$, with $\sum_gt_g=1$, such that $t\circ h$ is the same combination for all maps $h:[m]\to[n]$ with the same $m$.} For $n=1$, $M=2$, one can take $t=(0,1)-(0,2)+(1,2)$, naming a map by its values. In general, a map $g:[n]\to[M]$ is determined by its gaps $e_0=g(0)$, $e_i=g(i)-g(i-1)-1$, $e_{n+1}=M-g(n)$, that is, by a monomial $z^e$ of degree $M-n$ in $z_0,\dots,z_{n+1}$; a combination $t$ becomes a homogeneous polynomial $P_t$, and $\sum t_g=P_t(1,\dots,1)$. Let $D_i$ be the substitution $z_{i+1}\mapsto z_i$ followed by renumbering to variables $y_0,\dots,y_n$. Omitting the value $g(i)$ merges two gaps and the omitted point, so $P_{t\delta^i}=y_iD_i(P_t)$ for the map $\delta^i:[n-1]\to[n]$ that omits $i$. We construct, for all $n\geq0$ and $k\geq1$, a homogeneous polynomial $C_{n,k}$ with integer coefficients, $C_{n,k}(1,\dots,1)=1$, such that $y_i^kD_i(C_{n,k})$ does not depend on $i$. Put $C_{0,k}=1$. Given $Q=C_{n-1,2k}$ in the variables $y_0,\dots,y_n$, we want $D_i(C_{n,k})=F_i:=Q\prod_{l\neq i}y_l^k$ for all $i$; then $y_i^kD_i(C_{n,k})=Q\prod_ly_l^k$. In the coordinates $x_i=z_i-z_{i+1}$ ($0\leq i\leq n$), $u=z_{n+1}$, the map $D_i$ is reduction modulo $x_i$. The prescriptions agree modulo $(x_i,x_j)$: for $j=i+1$ both give $D_i(Q)\prod_lw_l^k$, and for $j\geq i+2$ they give $w_{j-1}^{2k}D_{j-1}(Q)\prod_{l\neq i,j-1}w_l^k$ and $w_i^{2k}D_i(Q)\prod_{l\neq i,j-1}w_l^k$, which are equal by the property of $Q$. For a nonempty set $S$ of indices let $\Phi_S$ be the common value of the $F_i$, $i\in S$, after setting $x_i=0$ for all $i\in S$ (in homogeneous lifts written in the coordinates $x,u$). Then
\[
 C_{n,k}=\sum_{S\neq\varnothing}(-1)^{|S|+1}\Phi_S
\]
reduces to $F_j$ modulo $x_j$, because the terms $S$ and $S\cup\{j\}$, $j\notin S$, cancel modulo $x_j$ and the term $S=\{j\}$ remains. It is homogeneous, and its value at $(1,\dots,1)$, a point where all $x_i$ vanish, is $F_0(1,\dots,1)=1$. Now let $t$ correspond to $C_{n,1}$. Then $s=t\delta^i$ is independent of $i$, and from $\delta^{b+1}\delta^a=\delta^a\delta^b$ ($a\leq b$) we get $s\delta^a=s\delta^b$; by induction on $n$, $s\circ h'$ depends only on the source of $h'$, and every non-identity $h:[m]\to[n]$ factors through some $\delta^i$. This proves the claim.

Let $F\subseteq F'$ be functors $\Delta'\to\mathrm{Ab}$. Any two objects of $\Delta'$ map to a third, so every element of $\varinjlim F$ is the class of one element $\xi\in F([n])$. Suppose its image in $\varinjlim F'$ is zero. The colimit is the quotient of $\bigoplus_mF'([m])$ by the elements $F'(v)\eta-\eta$, so $\xi=\sum_k(F'(v_k)\eta_k-\eta_k)$ in this direct sum, for finitely many $v_k:[s_k]\to[r_k]$ and $\eta_k\in F'([s_k])$. Choose $N$ and maps $d_m:[m]\to[N]$ for the finitely many $m$ that occur; applying $F'(d_m)$ on each summand gives $F'(d_n)\xi=\sum_k\bigl(F'(d_{r_k}v_k)-F'(d_{s_k})\bigr)\eta_k$ in $F'([N])$. Choose $t$ for $N$ and write $F'(t)=\sum_gt_gF'(g)$. Since $t\circ d_{r_k}v_k=t\circ d_{s_k}$, applying $F'(t)$ gives $F'(\tau)\xi=0$ for $\tau=t\circ d_n$, a combination of maps $[n]\to[M]$ with coefficient sum $1$. As $F\subseteq F'$, also $F(\tau)\xi=0$. The class of $F(g)\xi$ in $\varinjlim F$ is the class of $\xi$ for every $g$, so the class of $F(\tau)\xi$ is that of $\xi$; it is zero. So $\varinjlim_{\Delta'}$ preserves monomorphisms and is exact; $\Delta'$ is nonempty and connected, and the proposition applies.\hfill$\square$

That colimits of abelian groups over this category are exact is attributed to Isbell and Mitchell in \citet[A.5]{AdamekKoubekReiterman1981}, together with the general criterion: $\varinjlim_A$ is exact on abelian groups if and only if each component of the category with the same objects and the integral combinations of arrows of coefficient sum $1$ as morphisms is filtered. The withdrawn corollary on this page claimed $T$ for $\Delta^{\mathrm{op}}$, which is false; the category $\Delta'$ is an example of what that corollary was after.

\textbf{Theorem (monoids; finitely many objects).}\ours
\textit{Let $A$ have an object $q$ that receives an arrow from every object. Then $T(A)$ holds if and only if $A$ is filtered. This applies to every monoid, and to every category with finitely many objects, where both $T(A)$ and filteredness provide such an object $q$.}

\textit{Proof.}
Assume $T(A)$. \emph{Step 1: for parallel arrows $a,b:p\to q'$ there are an object $r$ and an integral combination $t$ of arrows $q'\to r$, with coefficient sum $1$, such that $ta=tb$.} Let $P_c=\mathbf Z[A(c,-)]$, let $d:P_{q'}\to P_p$ be $c\mapsto ca-cb$, with kernel $K$ and image $J$, and let $H$ be the quotient of $A(p,-)$ by the relations $ca\sim cb$, so that $0\to J\to P_p\to\mathbf Z[H]\to0$ is exact. The functors $P_c$ are projective with $\varinjlim P_c=\mathbf Z$, and $\varinjlim d=0$. So $L_1\varinjlim\mathbf Z[H]$ is the kernel of $\varinjlim J\to\varinjlim P_p$; this map is zero, because $\varinjlim P_{q'}\to\varinjlim J$ is surjective and $\varinjlim d=0$. Hence $L_1\varinjlim\mathbf Z[H]=\varinjlim J=\operatorname{coker}(\varinjlim K\to\mathbf Z)$, the image of $\varinjlim K$ being generated by the coefficient sums of the elements of the groups $K(r)$. The left side is $H_1(B\operatorname{El}(H))$, and $\operatorname{El}(H)$ is connected (every object receives an arrow from the class of $\mathrm{id}_p$); by the criterion it vanishes. So $1$ is a sum of coefficient sums of finitely many elements of groups $K(r_i)$; pushing them to a common successor $r$ of the $r_i$ (the common-successor proposition of the preceding commentary) and adding gives $t$.

\emph{Step 2: the monoid $M=A(q,q)$ is filtered.} For $x,m\in M$, Step~1 for the pair $x,xm$, followed by an arrow $r\to q$, gives $t\in\mathbf Z[M]$ with coefficient sum $1$ and $tx=txm$. The element $\nu=tx$ of $\mathbf Z[Mx]$ is invariant under right multiplication $\rho$ by $m$. Let $S$ be its finite support. Each point of $S$ has a $\rho$-preimage in $S$, distinct points have distinct preimages, so $\rho$ permutes $S$ and $\nu$ is constant on its cycles; the coefficient sum $1$ is thus an integral combination of cycle lengths. Among all points of $Mx$ that are $\rho$-periodic choose one, $w$, of least period $\pi$. If $\pi>1$, let $\ell$ be a prime dividing $\pi$ and apply the same to $w$ in place of $x$: some cycle in $Mw$ has length prime to $\ell$, and its points $w'$ satisfy $w'm^\pi=w'$ too, so their period is a proper divisor of $\pi$, a contradiction. Hence there is $y$ with $yxm=yx$. Given $m,m'\in M$, take $y$ with $ym=y$ (case $x=1$), then $y'$ with $y'ym'=y'y$; then $y'ym=y'y=y'ym'$.

\emph{Step 3.} Let $a,b:p\to q'$ be parallel, choose $g:q'\to q$, and by Step~1 an element $t=\sum n_ic_i\in\mathbf Z[M]$ of coefficient sum $1$ with $tga=tgb$. In the filtered monoid $M$ there is $y$ with $yc_i=y$ for all $i$. Then $yga=\sum n_iyc_iga=\sum n_iyc_igb=ygb$. With common successors this shows that $A$ is filtered. The converse is the fact that filtered categories satisfy $T$.\hfill$\square$

So Grothendieck's expectation is right in every finite case and for all monoids, and the smallest counterexamples have infinitely many objects. For the monoid $\{1,e,f\}$ with $uv=v$ for $u,v\in\{e,f\}$, which is cofiltered with contractible classifying space, the theorem says that $T$ fails; a cofinal image of the homotopy type of $S^2$ is obtained by adjoining a source $x$ with two arrows $h,i:x\to*$ on which $e$ and $f$ act by the constant map with value $i$. Whether $T(A)$ implies exactness of $\varinjlim_A$ on abelian groups, the converse of the proposition, is not decided here.
\end{dcomment}
